\documentclass{xduugthesis}
\usepackage{multirow}
\usepackage{amsmath}
\usepackage{float}
\usepackage{graphicx}
\usepackage{tabularx}
\usepackage{array}

\makeatletter
\let\cleardoublepage\clearpage
\makeatother
\setlength{\emergencystretch}{2em}
\setlength{\tabcolsep}{5pt}
\renewcommand{\arraystretch}{1.12}

\xdusetup{
  style / font-type = font,
  style / latin-font = tacn,
  style / cjk-font = fandol,
  info / title = {基于 AI Agent 与 MCP 的手机银行语义欺诈风控系统的设计与实现},
  info / author = {云天明},
  info / department = {计算机科学与技术学院},
  info / major = {软件工程},
  info / supervisor = {谢小璞、崔笛(西安电子科技大学)},
  info / class-id = {2203054},
  info / student-id = {22009200808},
  info / abstract = {abstract-zh.tex},
  info / abstract* = {abstract-en.tex},
  info / keywords = {AI Agent；MCP协议；手机银行；语义欺诈；风控系统},
  info / keywords* = {AI Agent; MCP Protocol; Mobile Banking; Semantic Fraud; Risk Control System},
  info / bib-resource = {references.bib},
  info / acknowledgements = {acknowledgements.tex}
}

\begin{document}

\chapter{绪论}

\section{研究背景与意义}

近年来，手机银行已经从辅助渠道逐渐变成用户办理转账、账户查询、理财和缴费等业务的主要入口。移动端交易的便利性提高了金融服务效率，但也改变了风险发生的位置：很多风险不再表现为设备被攻破、账号被盗用或异地登录，而是发生在用户本人主动操作的过程中。尤其在转账场景中，诈骗者往往通过冒充客服、虚假投资、公检法威胁等话术影响用户判断，使用户在设备、账号和网络环境都看似正常的情况下完成资金转出。这类风险具有较强的语义特征和行为隐蔽性，仅依靠传统规则很难被及时识别\cite{ahi2025risks}。

传统手机银行风控系统更擅长处理结构化异常，例如交易金额突增、登录地点变化、设备指纹异常或短时间高频交易等。这类规则对于盗刷和异常登录仍然有效，但面对语义欺诈时会遇到明显局限：交易动作由用户本人发起，交易参数也可能落在正常范围内，系统很难仅凭金额和设备信息判断用户是否正在受到诱导。也就是说，风险并不一定出现在交易字段本身，而可能隐藏在用户的交易意图、输入内容和操作节奏之中。

基于这一观察，本文将研究重点放在移动端转账链路中的隐性语义欺诈识别上。一方面，借助 AI Agent 对交易备注、用户补充说明和诈骗案例进行语义分析，尝试识别传统关键词规则不易覆盖的诱导性表达；另一方面，引入输入时长、停顿次数、修改频次和页面停留等微行为特征，补充描述用户在交易过程中的操作状态。由于金融交易对安全性和可追溯性要求较高，本文并不让 AI Agent 直接执行放行或拦截动作，而是通过规则治理层将语义分析结果收敛为确定性指令。这样的设计既利用了大模型在语义理解上的优势，也尽量保留金融风控系统对稳定性、权限边界和审计链路的要求。

因此，本文研究具有两方面意义。理论层面上，论文尝试把认知负荷、用户微行为和语义欺诈识别结合起来，为移动端风控提供一种不同于传统静态规则的分析视角。工程层面上，论文基于 Flutter、FastAPI、AI Agent 和 MCP 协议实现原型系统，验证“感知、认知、治理、执行”分层架构在手机银行转账场景中的可行性，为后续智能风控系统的安全落地提供参考。

\section{国内外研究现状}

从国外研究和产业实践来看，金融机构较早开始使用机器学习方法处理交易风险。随机森林、XGBoost 等模型常被用于识别大额异常交易、账户盗刷和支付欺诈，其优势在于能够从历史交易数据中学习结构化风险模式\cite{dubey2025transformer}。随着生成式人工智能的发展，一些研究开始关注大语言模型在金融文本处理、舆情分析、合规问答和交易描述理解中的应用。此类方法扩展了传统风控对非结构化文本的处理能力，但多数工作仍集中于通用金融文本分析，针对手机银行转账过程中“用户本人被诱导操作”的细粒度场景研究相对有限。

国内金融风控经历了从专家规则、评分卡模型到大数据用户画像的持续演进。近年来，围绕大模型轻量化部署、欺诈话术识别、多模态风险检测等方向已经出现较多探索\cite{nie2025multimodal}。这些研究能够在一定程度上识别诈骗文本或异常交易模式，但如果直接应用到手机银行核心交易场景，仍会遇到几个问题：一是模型输出具有概率性，难以直接作为资金交易的最终裁决；二是不同业务系统的数据接口不统一，影响风险上下文的完整获取；三是很多方案强调识别精度，却对审计追踪、权限隔离和人工复核等金融合规要求讨论不足。

AI Agent 技术为复杂任务拆解、上下文管理和工具调用提供了新的实现方式。在智能办公、问答系统和自动化流程中，Agent 已经能够根据任务目标调用外部工具并生成结构化结果。但在金融风控场景中，Agent 的角色需要受到更严格的限制。本文认为，Agent 更适合作为认知分析组件，而不应直接进入资金执行链路。原因在于大模型仍可能出现误判、幻觉或输出不稳定，如果缺少规则治理和审计约束，可能带来新的业务风险。

用户微行为分析是另一个与本文相关的研究方向。已有研究表明，用户在压力、诱导或高认知负荷状态下，输入节奏、停顿时间和修改行为可能发生变化。这类特征过去更多用于用户体验优化、身份识别或异常行为检测，在实时金融风控中的应用还不充分。本文将该方向与语义欺诈识别结合，尝试用移动端微行为补充交易文本和结构化字段的不足。

MCP 协议则为智能体与外部资源之间的标准化交互提供了新的接口思路。与传统自定义接口相比，MCP 更强调资源、工具和提示词的统一组织方式，便于在复杂系统中约束智能体的访问边界。现有应用多集中在通用智能体中台或开发者工具场景，面向银行强监管业务的定制化实践仍较少。基于此，本文将 MCP 引入手机银行风控原型系统，用于封装风险计算、转账执行和审计记录等工具，并配合权限隔离与 Trace ID 追踪机制，探索其在金融级审计链路中的应用价值。

\section{现有技术痛点分析}

结合手机银行转账业务特点和现有研究情况，本文认为当前移动端智能风控主要存在以下三个痛点。

首先，传统规则难以覆盖语义欺诈中的真实交易意图。现有风控系统通常根据交易金额、设备环境、登录地点和交易频次等字段进行判定，这些指标对于盗刷、撞库和异常登录有较好的识别效果。但是在社交工程类欺诈中，用户往往使用本人设备和本人账号完成转账，交易金额也可能被拆分到规则阈值以下。此时，风险并不表现为明显的参数异常，而是隐藏在“为什么转账”和“用户是否被诱导”之中。单纯依靠硬编码规则，很难捕捉这类语义层面的风险\cite{nicolaou2025llm}。

其次，风控上下文数据分散，影响系统对交易状态的整体判断。手机银行业务链路中可能同时涉及交易流水、设备信息、用户行为、风险案例库和历史处置记录等数据。如果这些数据分别存放在不同模块中，且缺少统一的调用规范，智能风控系统在分析时就难以获得完整上下文。对于 AI Agent 而言，上下文缺失会直接影响语义判断质量；对于审计流程而言，数据来源不清晰也会增加事后复核难度。因此，智能风控不仅需要识别模型，还需要一套规范的数据交互和追踪机制。

再次，大模型输出的不确定性与金融业务的确定性要求之间存在矛盾。大语言模型适合处理自然语言理解和复杂语义推理，但其输出并不是严格确定的。如果直接将模型给出的风险判断用于放行或拦截交易，一旦发生误判，就可能造成用户体验损害或资金安全风险。金融交易系统需要明确的权限边界、可解释的决策依据和可追溯的审计记录，因此必须将 AI 分析结果放入规则治理框架中进行二次收敛，而不能让模型直接控制业务执行。

基于上述问题，本文后续设计 HIRD 分层架构：感知层负责采集交易文本和微行为特征，认知层负责 Agent 语义分析，规则治理层负责将概率性结果转化为确定性指令，业务执行层负责完成放行、二次核验、拦截和人工复核等动作。该架构的出发点不是完全替代传统风控，而是在保留规则稳定性的基础上，引入语义理解和行为感知能力，补充传统手机银行风控在隐性语义欺诈场景中的不足。

\section{本文研究内容与组织结构}

\subsection{主要研究内容}

围绕上述问题，本文的研究内容主要集中在手机银行转账业务中的智能风控原型设计与实现。本文首先从交易链路出发，梳理正常转账、二次核验、风险拦截和人工复核等关键流程，明确系统需要处理的数据来源和权限边界。在此基础上，构建 HIRD 四层架构，将感知、认知、治理和执行拆分为相对独立的模块\cite{agnihotram2025agentic}。

在感知层，本文基于 Flutter 客户端实现微行为采集，重点记录输入时长、停顿次数、修改频次、页面停留和操作节奏方差等特征，并将其与交易金额、收款信息和交易备注一起封装为风险快照。在认知层，本文引入 AI Agent 和 RAG 检索机制，对交易文本、用户补充说明和风险案例进行语义匹配，输出风险线索和报告摘要。在规则治理层，本文设计动态权重与分级阈值机制，将 Agent 的概率性分析结果与硬规则评分进行整合，生成放行、二次核验或拦截等确定性指令。在业务执行层，本文通过 MCP 协议封装风险计算、转账执行和审计记录等工具，并保留 Trace ID、幂等校验和人工复核入口，使系统能够在测试环境中完成从感知采集到审计留痕的闭环。

\subsection{论文组织结构}

全文共分为六章，章节安排如下。

第一章为绪论，主要说明本文选题来源、研究背景、国内外研究现状和现有技术痛点，并给出本文的研究内容和整体组织方式。

第二章为相关技术综述，围绕本文系统实现所需的关键技术展开，介绍大语言模型与 AI Agent、MCP 协议、认知负荷理论、RAG 检索增强生成、Flutter 和 FastAPI 等内容，为后续架构设计提供技术基础。

第三章为系统需求分析与总体设计，结合手机银行转账场景拆解功能性和非功能性需求，进一步给出 HIRD 四层架构、微行为特征建模、Agent 协作流程、动态权重算法和审计链路设计。

第四章为系统实现与功能验证，说明前端、后端和 MCP 工具链的工程实现方式，并结合实际界面和接口流程展示系统如何完成风险预检、二次质询、报告生成和人工复核。

第五章为实验测试与分析，基于自动化测试脚本和场景化样本，对系统功能、响应时延、场景识别结果、影子模式一致性和消融实验结果进行分析。

第六章为结论与展望，总结本文完成的主要工作和实验结论，说明当前原型系统仍存在的数据规模、场景覆盖和消融验证不足，并提出后续改进方向。


\chapter{相关技术综述}
围绕本文构建的 HIRD（感知-认知-治理-执行）架构，本章介绍系统实现中实际用到的几类技术。与单纯罗列概念不同，本章更关注这些技术在手机银行风控原型中的作用边界：大语言模型与 AI Agent 负责语义分析，MCP 协议负责工具封装和调用约束，认知负荷理论为微行为特征提供解释依据，RAG、Flutter 和 FastAPI 则分别支撑知识检索、移动端采集和后端异步编排。
\section{大语言模型与金融AI Agent理论}
大语言模型通常基于 Transformer 架构，通过大规模语料训练获得较强的自然语言理解和生成能力\cite{nicolaou2025llm}。在手机银行风控场景中，它的价值主要体现在两点：一是能够处理交易备注、聊天片段、用户补充说明等非结构化文本；二是能够对同一类诈骗话术的变体进行语义层面的归纳，而不完全依赖固定关键词。例如，虚假理财、冒充客服和公检法威胁等场景中，诈骗者的具体措辞可能不断变化，但其诱导用户转账的意图往往具有相似结构。

不过，大语言模型并不适合直接作为金融交易的最终裁决器。模型输出可能受到提示词、上下文长度、案例覆盖程度和推理随机性的影响，存在误判或解释不稳定的风险。因此，本文在使用大模型时只把它放在认知分析环节，用于识别风险线索、生成语义摘要和辅助报告说明，而不让其直接执行放行、拦截或扣款等业务动作。

AI Agent 可以理解为在大语言模型基础上增加任务编排、上下文管理和工具调用能力的应用结构\cite{zangana2025llm}。与单次模型问答相比，Agent 更适合处理需要多步判断的任务，例如先读取交易上下文，再匹配风险案例，最后输出结构化风险原因。在本文系统中，Agent 主要位于 HIRD 架构的认知决策层，接收感知层生成的风险快照，并调用知识库检索、语义分析和报告生成等能力。

为了适配金融场景，本文对 Agent 的权限边界作了限制：Agent 可以分析、归纳和解释，但不能直接调用资金执行工具。系统最终的业务动作由规则治理层和业务执行层完成。这种设计使 Agent 能够发挥语义理解优势，同时避免模型输出直接影响交易状态。

\section{MCP深度解析}
Model Context Protocol（MCP）是一种面向智能体应用的上下文交互协议，主要用于规范模型、工具和外部资源之间的连接方式。对于本文系统而言，MCP 的作用不是提高模型本身的识别能力，而是为“Agent 能访问什么、能调用什么、调用后如何记录”提供统一约束。手机银行风控涉及交易数据、用户行为数据和审计日志，如果仍采用零散接口直接拼接，后续排查和权限管理都会变得困难。

MCP 协议中的 Resources、Tools 和 Prompts 三类机制分别对应不同问题。Resources 用于描述可访问的资源，例如风险案例库、交易上下文或审计记录；Tools 用于封装具体操作，例如风险评分、转账执行、PII 脱敏和日志落盘；Prompts 则用于约束 Agent 的输入输出格式，使语义分析结果更容易被后续规则治理层解析。通过这种方式，系统可以把复杂能力拆成边界清晰的原子工具，而不是让 Agent 直接接触底层业务逻辑。

在本文实现中，MCP 更接近一个安全数据中枢。Agent 需要通过 MCP 工具获取必要上下文，工具内部再完成权限检查、幂等校验和 Trace ID 记录。这样做的好处是，即便后续更换模型或调整提示词，底层交易执行和审计链路仍由固定工具控制。对金融风控系统来说，这种可追踪、可复核的调用方式比单纯追求模型输出更重要。

\section{认知负荷理论与微行为特征分析}
认知负荷理论认为，人在处理信息时会受到注意力、短时记忆和任务复杂度的限制。当用户在转账过程中同时面对诈骗者话术、时间压力、身份伪装或威胁暗示时，其认知资源会被额外占用，操作行为也可能出现变化\cite{ibraheem2025duress}。这类变化不一定会体现在交易金额或设备信息中，但可能反映在输入速度、停顿位置、修改次数和页面停留时间等细节上。因此，在语义欺诈识别中，仅分析交易字段是不够的，还需要观察用户完成交易动作时的过程特征。

本文没有把微行为特征直接等同于欺诈结论，而是将其作为辅助风险证据使用。原因在于，用户输入变慢或频繁修改并不一定意味着被诈骗，也可能来自网络卡顿、年龄差异、输入习惯或临时犹豫。因此，本文在感知层采集有效输入时长、非预期停顿次数、内容修改频次、页面停留时长和操作节奏方差等指标后，并不单独据此拦截交易，而是将其与交易文本、设备状态和规则评分共同送入后续风控链路。

从系统设计角度看，微行为特征的价值在于补充传统结构化字段的盲区。对于“用户本人主动转账”的社交工程欺诈，设备和账号可能完全正常，金额也可能被控制在规则阈值以内。此时，用户在交易界面上的停顿、反复修改和异常节奏，能够为认知决策层提供额外线索。本文后续在 HIRD 架构中将这类数据封装为感知快照，使 Agent 和规则治理层能够在同一条风险链路中使用这些信息。

需要注意的是，微行为特征存在较强的个体差异，不能简单套用统一阈值。本文原型系统主要通过场景化样本验证该思路的可行性，后续若接入真实业务环境，还需要结合用户历史行为基线、设备性能差异和更多异常样本进行校准。

\section{检索增强生成(RAG)与知识库构建}
检索增强生成（RAG）的基本思路，是在大语言模型生成回答前，先从外部知识库中检索与当前问题相关的材料，再把检索结果作为上下文提供给模型\cite{srinivasan2025rag}。在本文场景中，RAG 的作用不是让模型“知道更多金融知识”这么简单，而是让 Agent 的风险判断能够尽量依托具体案例和规则片段，减少仅凭模型内部知识进行推断带来的不稳定性。

语义欺诈具有明显的变体特征。同一类诈骗可能使用不同表达，例如虚假理财可以表现为“内部收益通道”“限时认购”“稳赚回款”等多种话术；冒充客服也可能通过订单异常、退款验证、账户冻结等不同理由诱导转账。单纯依靠关键词匹配时，规则库需要不断枚举具体词语，维护成本较高。RAG 结合向量检索后，可以从语义相似度角度匹配历史案例，使系统对同类变体话术具备一定的泛化能力。

本文知识库主要围绕三类内容组织：一是典型欺诈案例，用于提供场景匹配依据；二是风控规则说明，用于约束 Agent 输出的风险类型和处置建议；三是系统报告模板，用于规范风险摘要和风险因子字段。Agent 在处理交易备注或用户补充说明时，会先检索相似案例，再结合当前交易金额、收款信息和微行为快照进行判断。这样生成的风险说明更容易追溯到具体案例，也方便后续人工复核。

同时，RAG 也存在局限。检索结果的质量依赖知识库覆盖范围和向量匹配效果，如果案例库过小或样本过于模板化，Agent 仍可能遗漏新型诈骗话术。本文当前知识库主要服务于原型系统和场景化测试，后续需要通过更丰富的真实或半真实案例持续补充，才能支撑更复杂的生产环境。

\section{跨平台开发与异步编排技术}
本文原型系统选择 Flutter 作为移动端实现框架，主要考虑其跨平台能力和事件监听机制。手机银行风控采集不能明显打断用户操作，也不能让用户为了测试而额外输入数据。因此，感知层需要在转账流程中记录输入时长、停顿次数、修改频次和页面停留等行为数据，并在用户提交交易时生成统一快照。Flutter 的组件化界面和事件回调机制适合承载这种采集逻辑，也便于在同一套代码中模拟不同移动端交互流程。

在本文实现中，移动端采集不是为了记录用户隐私内容，而是为了形成可计算的行为指标。系统关注的是“输入用了多久”“是否出现非预期停顿”“是否反复修改”等过程信息，而不是简单保存用户完整输入过程。这样做可以在保留风险判断线索的同时，减少不必要的敏感数据暴露。后续与 MCP 数据中枢结合时，这些指标会和交易文本、交易元数据一起进入风险快照。

后端部分采用 FastAPI，主要原因是风控链路中存在多类耗时差异明显的任务：预检热路径需要在 100ms 左右完成，二次核验和 Agent 语义分析则可能受到大模型接口、网络请求和检索过程影响。FastAPI 基于异步机制组织接口调用，便于把本地规则判断、日志记录、模型调用和状态机流转拆开处理。在第五章测试中，预检路径和二次核验路径的时延差异也说明，系统需要通过双路径方式处理不同复杂度的交易请求。

因此，Flutter 和 FastAPI 在本文中并不是单纯的前后端技术选型，而是分别服务于“无感采集”和“异步风控编排”两个目标。前者负责在移动端捕捉交易过程，后者负责在后端把规则、Agent、MCP 工具和审计日志组织为可追踪流程。

\section{本章小结}
本章围绕 HIRD 架构所需的关键技术进行了说明。大语言模型和 AI Agent 为交易文本、用户说明和欺诈案例之间的语义分析提供能力，但在本文系统中只承担认知分析职责，不直接控制业务执行。MCP 协议用于封装资源和工具调用，使 Agent 的访问边界、执行动作和审计记录能够被统一约束。认知负荷理论和微行为特征分析为移动端操作过程提供解释依据，使系统能够从用户输入节奏和停顿行为中获得补充风险线索。

在工程实现方面，RAG 与知识库用于增强 Agent 对欺诈案例的检索和引用能力，降低单纯依赖模型内部知识带来的不稳定性；Flutter 支撑移动端无感采集，FastAPI 支撑后端异步编排和双路径决策。上述技术共同服务于本文的核心目标：在不让 AI 直接接管资金执行的前提下，将语义理解、行为感知和规则治理纳入同一条可审计的手机银行风控链路。后续第三章将在这些技术基础上展开系统需求分析和 HIRD 架构设计。

\chapter{系统需求分析与总体设计}

前两章已经说明，本文关注的不是一般意义上的交易异常检测，而是手机银行转账过程中较难被传统规则识别的语义欺诈。围绕这一场景，本章从实际业务链路出发，先拆解系统需要满足的功能与约束，再给出 HIRD 分层架构、感知特征建模和后续决策链路设计。章节中的需求并非独立罗列，而是服务于同一目标：在不让 AI 直接控制资金执行的前提下，把移动端行为、交易语义和规则治理组织成可测试、可审计的风控流程。

\section{系统业务需求分析}

本文原型系统以手机银行转账为主要业务对象。一次完整交易通常包含用户填写信息、系统预检、必要时触发二次核验、最终放行或拦截，以及后续查看风险报告和人工复核等步骤。因此，需求分析需要同时覆盖前端采集、后端判断、AI 分析、规则收敛和审计留痕，而不能只关注某一个识别模型的准确率。

\subsection{功能性需求}

结合本文系统实现，功能性需求被划分为五类：多模态风险感知、语义欺诈识别、安全数据交互、动态分级治理和全链路审计。前两类需求负责发现风险，中间两类需求负责约束 AI 分析和生成确定性处置，最后一类需求用于支撑事后复核。

\textbf{F-01 多模态风险感知功能}

在社交工程类欺诈中，用户往往使用本人设备和本人账号发起交易，金额、IP 地址等字段也可能没有明显异常。为避免系统只看到“参数正常”的一面，感知模块需要在不打断用户转账流程的前提下记录操作过程，例如停顿、修改、页面停留和设备上下文，并将这些信息与交易参数放入同一份风险快照中。

\begin{table}[H]
    \centering
    \caption{多模态风险感知功能细分}
    \label{tab:f01-sub}
    \begin{tabular}{|c|p{7cm}|}
        \hline
        细分功能名称 & 核心任务说明 \\
        \hline
        异步非阻塞行为监听 & 利用前端框架在后台静默采集按键停顿等数据，不干扰转账界面正常交互 \\
        \hline
        微行为数据量化 & 将输入停顿时间、反复修改动作等操作特征转换为系统可计算的数值指标 \\
        \hline
        交易环境特征整合 & 同步提取登录位置、设备型号等基础结构化信息，与行为数据对齐 \\
        \hline
        风险感知快照生成 & 将多源数据按时间线对齐，封装为固定格式的快照供下游模块消费 \\
        \hline
    \end{tabular}
\end{table}

\textbf{F-02 语义欺诈智能识别功能}

很多诈骗话术并不会稳定使用同一组关键词。即使都属于虚假理财或冒充客服场景，用户看到的表达也可能不断变化。因此，系统需要能够分析交易备注、用户补充说明和案例库中的语义关系，识别隐藏在自然语言背后的诱导意图。该功能输出的结果只作为风险线索和评分输入，不直接转化为业务执行动作。

\begin{table}[H]
    \centering
    \caption{语义欺诈智能识别功能细分}
    \label{tab:f02-sub}
    \begin{tabular}{|c|p{7cm}|}
        \hline
        细分功能名称 & 核心任务说明 \\
        \hline
        复杂话术意图提取 & 解析转账附言或聊天记录中的诱导性文本，识别其真实交易目的 \\
        \hline
        知识库相似度比对 & 将文本特征与RAG知识库中的历史欺诈案例进行语义级相似度匹配 \\
        \hline
        隐性风险分值计算 & 综合上下文语义理解与案例匹配结果，输出量化的语义风险评分 \\
        \hline
    \end{tabular}
\end{table}

\textbf{F-03 安全数据交互与权限管控功能}

Agent 在分析交易风险时需要读取上下文，但银行业务数据具有较高敏感性，不能简单把原始交易信息直接暴露给模型。本功能要求系统在接口层面限制 Agent 的访问范围，只提供完成语义分析所需的脱敏特征和必要上下文，并通过 MCP 工具封装方式记录每一次调用。

\begin{table}[H]
    \centering
    \caption{安全数据交互与权限管控功能细分}
    \label{tab:f03-sub}
    \begin{tabular}{|c|p{7cm}|}
        \hline
        细分功能名称 & 核心任务说明 \\
        \hline
        数据访问权限隔离 & 限制大模型仅能读取脱敏后的特征数据，禁止直接访问底层数据库 \\
        \hline
        原子化风控工具封装 & 基于MCP协议将风险计算等动作封装为独立工具，规范调用路径 \\
        \hline
        敏感隐私信息脱敏 & 在数据传递给外部模型前，自动对姓名、卡号等字段进行掩码处理 \\
        \hline
    \end{tabular}
\end{table}

\textbf{F-04 动态分级风控决策与治理功能}

大模型输出更适合描述风险可能性，而金融交易需要明确动作。因此，系统需要在治理层把语义评分、规则评分和场景匹配程度重新整合，输出放行、二次核验或拦截等确定性结果。该功能的重点不是让 AI 取代规则，而是在规则框架内吸收语义分析的补充价值。

\begin{table}[H]
    \centering
    \caption{动态分级风控决策与治理功能细分}
    \label{tab:f04-sub}
    \begin{tabular}{|c|p{7cm}|}
        \hline
        细分功能名称 & 核心任务说明 \\
        \hline
        场景匹配度计算 & 评估当前交易与已知诈骗模式的相似程度，决定两类评分的权重分配 \\
        \hline
        动态权重参数调节 & 当交易呈现异常特征时，自动提升AI评分的权重占比进行融合计算 \\
        \hline
        综合风险等级划分 & 依据融合后的总分值，将交易判定为高、中、低三个风险等级 \\
        \hline
        分级路由指令下发 & 根据风险等级自动执行放行、人脸核验或交易拦截等差异化处置 \\
        \hline
    \end{tabular}
\end{table}

\textbf{F-05 全链路风控审计溯源功能}

对于被拦截或进入二次核验的交易，系统需要保留足够的过程信息，方便用户查看、后台复核和后续问题定位。本功能要求将感知快照、语义分析结果、权重计算过程、最终指令和人工复核状态纳入同一条审计链路，避免只保存一个简单的“拦截”结果。

\begin{table}[H]
    \centering
    \caption{全链路风控审计溯源功能细分}
    \label{tab:f05-sub}
    \begin{tabular}{|c|p{7cm}|}
        \hline
        细分功能名称 & 核心任务说明 \\
        \hline
        原始输入快照存档 & 完整保存系统采集的行为参数与设备信息，作为事后核查的底层依据 \\
        \hline
        决策推导过程记录 & 将大模型话术分析、案例匹配情况及权重分配逻辑写入审计日志 \\
        \hline
        处置指令存储 & 将最终指令及操作时间持久化存储，确保记录不可篡改 \\
        \hline
    \end{tabular}
\end{table}

\subsection{非功能性需求}

非功能性需求决定了系统能否真正放入转账链路中运行。本文主要从实时性、安全性、准确性和可扩展性四个方面给出约束指标，如表\ref{tab:nonfunctional-req}所示。这些指标后续将在第五章通过自动化测试、场景样本和性能统计进行验证。

\begin{table}[H]
    \centering
    \caption{系统非功能性需求}
    \label{tab:nonfunctional-req}
    \begin{tabular}{|c|c|p{5cm}|}
        \hline
        需求编号 & 指标项 & 具体要求 \\
        \hline
        NF-01 & 实时性 & 预检≤100ms，含AI链路≤3s \\
        \hline
        NF-02 & 安全性 & 敏感信息全脱敏，接口幂等校验 \\
        \hline
        NF-03 & 准确性 & 召回率≥90\%，误报率≤10\% \\
        \hline
        NF-04 & 可扩展性 & 层级解耦，可独立升级与扩充 \\
        \hline
    \end{tabular}
\end{table}

\textbf{NF-01 实时性}

转账预检处在用户提交交易后的关键节点，如果等待时间过长，会直接影响手机银行的交互体验。因此，本文将预检热路径控制在 100ms 以内，该路径只执行本地规则、结构化字段判断和必要的状态记录。当交易进入二次核验或 Agent 语义分析时，系统允许更长的异步处理时间，并通过前端等待状态和后端状态机避免阻塞主流程。

\textbf{NF-02 安全性}

安全性主要体现在三个方面：敏感信息在进入 Agent 分析前需要脱敏；接口调用需要进行权限校验；涉及交易执行或状态变更的 MCP 工具需要具备幂等保护。这样即使出现网络重试、重复提交或并发调用，也不会造成重复扣款、重复审计或交易状态漂移。

\textbf{NF-03 准确性}

准确性需要同时考虑欺诈召回和正常交易误报。本文希望系统能够尽量识别典型语义欺诈样本，同时避免对正常交易造成过度干预。由于当前实验数据主要为场景化样本，本文将该指标作为原型系统测试目标，并在第五章中结合 54 个核心场景样本进行统计分析。

\textbf{NF-04 可扩展性}

可扩展性主要来自分层解耦。感知层可以增加新的行为指标，认知层可以替换模型或扩充知识库，治理层可以调整权重和阈值，执行层可以增加新的复核或提示方式。各层通过标准化快照、MCP 工具和明确状态字段衔接，尽量减少局部修改对整体系统的影响。

综合来看，上述约束之间存在明显取舍。实时性要求系统不能在每一笔交易中都调用大模型；识别能力又要求系统在规则难以判断时引入语义分析；安全性则要求 Agent 不能直接访问底层业务数据或执行交易动作。因此，本文采用热路径与语义路径分流的方式：常规交易优先由本地规则快速处理，疑似风险交易再进入二次核验和 Agent 分析；同时通过 MCP 工具和规则治理层把分析权限与执行权限分开。这一取舍构成了 HIRD 架构的设计出发点。

\section{HIRD 系统总体架构设计}

根据 3.1 节的需求拆解，本文提出 HIRD 四层分层风控架构。HIRD 分别对应感知集成层、认知决策层、规则治理层和业务执行层。该架构的核心思路是让不同能力停留在各自边界内：感知层只负责形成数据，认知层只负责分析风险，治理层负责把分析结果转成明确指令，执行层只执行已经确定的业务动作。这样设计的目的，是避免 AI 分析结果绕过规则治理直接影响资金流转。

\subsection{架构层级定义与设计理念}

本文将系统拆解为四个层级，并不是为了增加架构复杂度，而是为了明确每一类能力的边界。感知集成层（H层）负责采集交易文本、结构化参数和微行为数据；认知决策层（I层）负责语义分析、案例匹配和风险线索归纳；规则治理层（R层）负责融合规则评分和语义评分，输出确定性处置；业务执行层（D层）负责前端提示、交易放行、拦截、报告展示和人工复核入口。架构总览如图\ref{fig:hird-architecture}所示。

% ===== 手动插图占位（已禁用，避免编译报错）=====
% 原图路径：figures/ui/candidate_p31_2_Im5.png
% \begin{figure}[h]
%     \centering
%     \includegraphics[width=0.82\linewidth]{figures/ui/candidate_p31_2_Im5.png}
%     \caption{HIRD四层风控架构图}
%     \label{fig:hird-architecture}
% \end{figure}

上述四层划分的关键，是把“分析权限”和“执行权限”分开。I层可以读取脱敏后的上下文并生成风险线索，但不能直接调用转账执行或拦截工具；R层根据规则和阈值决定最终指令；D层只接收明确指令并完成页面反馈或状态变更。这样即使 Agent 的语义判断出现偏差，也不会直接越过治理层影响资金动作。

\subsection{层间数据流转与协作机制}

一笔交易从用户提交到最终完成，会在 HIRD 四层之间按固定方向流转。系统使用全局唯一交易流水号贯穿感知快照、语义分析、治理决策和审计记录，确保每个环节都能被定位。根据预检结果不同，交易可能走快速放行或拦截路径，也可能进入二次质询、报告生成和人工复核路径。
% ===== 手动插图占位（已禁用，避免编译报错）=====
% 原图路径：figures/ui/candidate_p31_1_Im4.png
% \begin{figure}[h]
%     \centering
%     \includegraphics[width=0.92\linewidth]{figures/ui/candidate_p31_1_Im4.png}
%     \caption{各层间数据流转}
%     \label{fig:hird-dataflow}
% \end{figure}

\textbf{预检路径}覆盖所有交易的首轮判定。用户提交转账请求后，H层生成包含交易参数、设备上下文、文本内容和微行为指标的感知快照。快照进入I层后，系统对交易文本进行基础分类和案例匹配，形成风险类别、匹配场景和语义线索。R层再结合结构化规则、行为特征和语义结果计算综合风险分值，并映射为pass、block或interrogate三类结果。pass表示交易可继续执行，block表示直接拦截，interrogate表示需要进入二次质询。整个预检过程的关键节点写入审计日志，前端同步展示决策结果和必要解释信息。

\textbf{二次质询路径}用于处理预检阶段难以直接放行或拦截的交易。R层先检查交易状态是否允许进入二次质询，随后由前端提示用户补充说明。I层的 Agent 服务会先根据本地红旗规则判断是否存在明确恶意回复，再判断是否属于低风险合规回复；当本地规则无法覆盖时，再调用大模型进行语义分析。分析结果回到R层后，R层更新交易状态、写入审计日志，并触发风险报告生成。前端随后获得二次决策结果、语义红旗标识和解释面板数据。

\textbf{风险报告查看与人工复核路径}用于补充自动化风控的解释和兜底能力。用户可以在前端查看风险标题、综合风险等级、风险摘要和风险因子。如果用户对拦截或核验结果有异议，可以从报告详情页发起复核申请。系统创建复核记录后进入后台处理队列，审核人员可查看交易上下文、审计链路和系统解释，再给出处置结论。复核结果作为治理记录留存，不直接回写自动化裁决，以避免人工流程破坏原有决策链路的可追溯性。

\textbf{确认与取消路径}是交易状态的最终出口。确认操作只允许在交易状态满足条件时执行；如果交易经过二次质询，则必须先获得通过结果。取消操作允许用户在待处理状态下主动终止交易，取消后系统拒绝后续确认或复核外的状态变更。确认和取消都会写入审计日志，用于还原交易生命周期。

\begin{table}[H]
    \centering
    \caption{各层级在不同流转路径中的职责}
    \label{tab:layer-responsibility}
    \begin{tabular}{|c|p{3cm}|p{3cm}|p{3cm}|}
        \hline
        层级 & 预检路径 & 二次质询路径 & 报告与复核路径 \\
        \hline
        H层（感知） & 聚合多源数据，封装感知快照 &  &  \\
        \hline
        I层（认知） & 文本分类，产出风险标签 & 语义分析，产出决策建议 &  \\
        \hline
        R层（规则治理） & 公式计算，阈值路由 & 状态更新，触发报告生成 & 持久化报告，管理复核状态机 \\
        \hline
        D层（执行） & 执行放行/拦截 & 展示解释面板，执行裁决 & 提供报告入口，展示时间线 \\
        \hline
    \end{tabular}
\end{table}

\section{感知集成层（H）特征建模}

感知集成层是 HIRD 架构的入口。它不负责判断交易是否欺诈，而是负责把一次转账过程中分散出现的文本、结构化参数和微行为数据整理成后续模块可以使用的快照。对于本文系统来说，H层的设计重点有两个：一是采集过程不能明显打断用户操作；二是采集结果必须足够规范，能够被 Agent、规则引擎和审计模块共同引用。

\subsection{微行为特征指标体系}

结合第二章对认知负荷和微行为的分析，本文在移动端转账流程中选取五项较容易采集、也较容易解释的行为指标。选择这些指标的原因是，它们都与用户完成转账动作的过程有关，而不是只描述交易最终结果。例如，输入时间变长、停顿增多、反复修改或确认前长时间停留，都可能提示用户正在受到外部信息干扰。当然，这些指标不能单独证明欺诈，只能作为后续综合判断的辅助证据。各指标的判异规则如表\ref{tab:micro-behavior}所示。

\begin{table}[H]
    \centering
    \caption{微行为特征指标及判异规则}
    \label{tab:micro-behavior}
    \begin{tabular}{|c|p{8cm}|}
        \hline
        指标 & 判异规则 \\
        \hline
        有效输入时长 & 剔除外部干扰后的填写耗时；若较用户基线显著延长则判定异常 \\
        \hline
        非预期停顿次数 & 操作间隔超过500ms的停顿数；若次数增多且分布杂乱则判定异常 \\
        \hline
        内容修改频次 & 关键字段的修改次数；频繁修改反映外部指令引导 \\
        \hline
        页面停留时长 & 填完到确认的驻留时间；长时间停留反映决策犹豫 \\
        \hline
        操作节奏方差 & 相邻操作间隔的方差；方差增大反映注意力分散和节奏紊乱 \\
        \hline
    \end{tabular}
\end{table}

以上五项指标在系统中以特征向量形式使用。本文并不为某一个指标设置绝对拦截条件，而是更关注多项指标的组合偏离。例如，单次页面停留时间变长可能只是用户临时犹豫，但如果同时出现输入时间延长、修改次数增加和操作节奏方差变大，就更值得进入后续语义分析或二次核验。这样的处理方式能够降低单一行为指标带来的误判风险。

\subsection{标准化感知快照生成机制}

移动端采集到的原始事件通常比较碎片化，不能直接交给后端模型或规则引擎使用。H层需要先对事件进行整理，再形成统一格式的感知快照。本文将这一过程划分为噪声过滤、归一化和快照封装三个步骤。

\textbf{噪声过滤}主要处理设备卡顿、网络延迟、误触和页面渲染造成的异常事件。例如，毫秒级间隔内的连续重复点击更可能来自设备抖动，不应直接计入停顿次数；页面加载造成的等待也需要与用户真实停留区分开。过滤步骤的目标不是追求行为记录越多越好，而是尽量保留与用户决策过程相关的数据。

\textbf{归一化}用于处理不同指标的量纲差异。输入时长通常以秒计，修改频次是离散次数，操作节奏方差则可能落在更小的数值范围内。为避免某一类数值因为量纲较大而在后续计算中占据过高权重，系统采用 Min-Max 方法将连续型指标映射到[0,1]区间。

\textbf{快照封装}负责把处理后的行为向量、交易结构化数据和语义文本数据绑定到同一笔交易流水号下，形成标准化感知快照。每笔交易对应一条主要快照，后续I层、R层和审计模块都围绕该快照继续处理。这样可以减少跨层传递时的数据混淆，也方便在风险报告和人工复核中回溯原始依据。

\section{认知决策层（I）详细设计}

认知决策层位于感知快照和规则治理之间，负责把“已经整理好的交易上下文”转化为可供治理层使用的风险线索。它处理的对象包括交易备注、用户补充说明、风险案例匹配结果和部分脱敏后的行为特征。I层的定位仍然是分析层：它可以给出语义红旗、风险摘要和案例匹配依据，但不直接决定交易放行或拦截。这样的边界设置，能够让 Agent 参与复杂语义判断，同时避免模型输出直接进入资金执行链路。

\subsection{多Agent协作架构}

本文将I层拆分为四类 Agent，主要是为了降低单个 Agent 同时承担分类、推理、解释和报告生成时带来的不稳定性。初分 Agent 负责快速判断交易属于哪类场景；细分 Agent 在场景范围内匹配更具体的欺诈模式；语义分析 Agent 关注文本中是否存在诱导、威胁、收益承诺等风险表达；报告生成 Agent 则把前面几类结果整理成用户和审核人员能够阅读的结构化说明。各 Agent 的角色定义与输出约束如表\ref{tab:agent-constraint}所示。

\begin{table}[H]
    \centering
    \caption{Agent角色描述与约束矩阵}
    \label{tab:agent-constraint}
    \begin{tabular}{|c|p{5cm}|p{5cm}|}
        \hline
        Agent角色 & 核心任务 & 输出约束 \\
        \hline
        初分Agent & 对交易场景进行快速归类 & 仅输出场景类别标签，不做深度判定 \\
        \hline
        细分Agent & 在初分类别下识别具体欺诈模式 & 输出风险评分及匹配案例编号 \\
        \hline
        语义分析Agent & 解析交易文本中的欺诈意图 & 输出语义红旗标识与推理依据摘要 \\
        \hline
        报告生成Agent & 将推理过程整理为可读报告 & 仅做说明性输出，不参与业务裁决 \\
        \hline
    \end{tabular}
\end{table}

四类 Agent 的共同约束是只输出分析结果，不直接调用业务执行工具。前三类 Agent 的输出主要服务于规则治理层，报告生成 Agent 的输出主要服务于解释面板和审计报告。这样设计后，即使某一类 Agent 给出偏高或偏低的判断，也必须经过 R 层的权重、阈值和状态校验后，才能影响最终处置。

\subsection{Agent协作流程}

四类 Agent 的协作遵循从粗到细的顺序。系统先判断交易大类，再匹配具体风险模式，随后提取语义红旗，最后生成说明性报告。这样安排的好处是，后续 Agent 不需要面对完全开放的文本空间，而是可以在前一阶段已经缩小的场景范围内继续分析。

感知快照从H层进入I层后，首先由初分 Agent 接收。初分 Agent 依据交易金额、收款信息、设备上下文和文本初步特征，将当前交易归入常规转账、疑似虚假理财、疑似冒充客服、疑似公检法诱导等场景。该阶段只输出类别标签和置信度，不给出最终风险结论。

初分结果产生后，细分 Agent 会在对应类别下检索案例库，并结合外部情报判断当前交易是否接近已知风险模式。例如，若初分结果指向虚假理财，系统会优先匹配收益承诺、限时认购、账户归集等案例特征；若指向冒充客服，则会关注退款验证、订单异常和账户冻结等诱导理由。细分 Agent 输出相似度评分和案例编号，用于保留判断来源。

在细分匹配完成后，语义分析 Agent 进一步处理交易备注、用户补充说明等非结构化文本。该阶段不再只看场景类别，而是提取具体风险表达，例如要求转入“安全账户”、承诺固定收益、强调保密或制造紧迫感等。分析结果以语义红旗列表输出，每个红旗附带简短解释，说明它为什么可能提升交易风险。

前三步完成后，报告生成 Agent 汇总场景类别、案例匹配结果和语义红旗，生成结构化风险说明。报告内容包括风险标题、摘要、风险因子和必要的解释文本。该报告主要用于前端展示、后台复核和审计查询，不参与最终裁决；最终处置仍由 R 层根据评分和阈值完成。


\section{规则治理层（R）决策算法设计}

规则治理层负责把I层给出的风险线索转化为可执行的业务指令。由于 Agent 输出带有概率性和解释性，不能直接决定交易状态，R层需要在硬规则、语义评分和外部风险情报之间做一次收敛。本节重点说明两个问题：一是如何控制 AI 语义评分在最终结果中的权重；二是如何把连续风险分值映射到放行、二次核验和拦截三类动作。

\subsection{动态权重调整机制}

传统规则引擎的优势是稳定、可解释，适合处理金额阈值、设备风险、交易频次等明确条件；不足是很难覆盖不断变化的诈骗话术。AI 语义分析能够识别更灵活的表达，但也可能受到提示词、案例覆盖和模型波动影响\cite{haeri2025riskembed}。因此，本文没有让其中任意一方单独决定结果，而是通过权重机制把两类信号放到同一个评分框架中。

R层在收到感知快照和I层分析结果后，会先查看当前交易与知识库案例的相似度，再结合设备、账户和交易参数等结构化风险信息，决定 AI 语义评分的权重调节因子 $\omega_{adj}$。该参数并不是由 Agent 自己决定，而是由治理层根据规则计算得到，用于限制模型输出对最终分值的影响。

\begin{itemize}
    \item 当交易与已知欺诈模板相似度较低，且外部情报没有明显异常时，$\omega_{adj}$取0.1。此时最终结果主要由硬规则决定，AI 语义结果只作为弱参考，避免正常交易被模型误判影响。
    \item 当交易匹配到高风险案例，或设备、账户等外部情报提示异常时，$\omega_{adj}$取0.7。此时系统提高语义评分影响力，用于补充规则对变体话术识别不足的问题。
\end{itemize}

\subsection{最终风险评分模型}

在完成权重设置后，R层计算最终综合风险分值。该模型保留两部分输入：一部分来自I层输出的语义风险得分，另一部分来自传统结构化风险特征。公式如下：

\begin{equation}
f_{final} = \omega_{adj} \times R_{agent} + \sum_{i=1}^{n} \omega_i \times x_i
\label{eq:final-score}
\end{equation}

式中各参数含义如下：

\begin{itemize}
    \item $f_{final}$：最终综合风险得分，取值范围[0,100]，分数越高代表交易风险等级越高；
    \item $\omega_{adj}$：AI权重调节因子，取值范围[0,1]，由R层规则动态生成，用于控制语义风险得分对最终结果的影响；
    \item $R_{agent}$：认知层输出的语义风险得分，取值范围[0,100]，表征交易文本与用户行为的综合语义风险程度；
    \item $\omega_i$：第$i$项传统结构化特征静态权重，由系统配置给出；
    \item $x_i$：第$i$项结构化风险特征数值，包含交易金额、设备风险等级、交易地域等参数。
\end{itemize}

该模型的设计重点在于约束，而不是单纯叠加分数。当交易缺少明显语义风险时，$\omega_{adj}$较低，系统仍以结构化规则为主；当交易与欺诈案例高度相似或外部情报异常时，语义评分的影响增大。这样可以在不放弃规则稳定性的前提下，引入 Agent 对复杂文本的分析能力。

\subsection{分级治理与人工复核路由}
得到 $f_{final}$ 后，系统需要把连续分值转化为用户能够感知的业务动作。本文采用三级分级方式：低风险尽量不打扰用户，中风险通过二次核验进一步确认，高风险直接阻断资金流转。分级映射如表\ref{tab:risk-grading}所示。

\begin{table}[H]
    \centering
    \caption{风险分级处置映射表}
    \label{tab:risk-grading}
    \begin{tabular}{|c|c|c|c|}
        \hline
        风险分值区间 & 风险等级 & 处置动作 & 用户侧表现 \\
        \hline
        $0 \le \text{得分} < 30$ & 低风险 & 自动放行 & 无感知，交易正常完成 \\
        \hline
        $30 \le \text{得分} < 70$ & 中风险 & 二次核验 & 弹出质询弹窗或人脸验证 \\
        \hline
        $70 \le \text{得分} \le 100$ & 高风险 & 直接拦截 & 阻断交易并展示风险提示 \\
        \hline
    \end{tabular}
\end{table}

低风险区间表示当前交易缺少明显异常，系统不额外打扰用户。中风险区间表示系统发现一定风险线索，但证据不足以直接拦截，此时触发二次质询、人脸核验或补充说明；若交易金额较大或分值接近高风险边界，还可以进入人工复核队列。高风险区间表示交易与已知欺诈模式或强风险特征高度吻合，系统直接拦截并生成风险报告，同时保留完整审计数据供后续核查。


\section{业务执行层（D）与全链路审计设计}

业务执行层是 HIRD 架构中最接近用户和交易状态的一层。前三层完成数据采集、语义分析和规则收敛后，D层只接收R层已经确定的指令，并把它转化为页面提示、交易放行、交易拦截、报告展示或复核入口等具体动作。本文将D层设计为“只执行、不分析”，目的是避免执行层再次引入模糊判断，也避免 Agent 绕过治理层直接影响交易状态。

D层的设计有两个重点。第一是状态确定性：系统必须清楚一笔交易当前处于待确认、二次核验、已放行、已拦截、已取消或复核中等哪一种状态。第二是过程可追溯：每一次状态变化都要记录触发来源、时间、输入数据和执行结果。基于这两点，D层主要承担指令执行、用户交互、MCP工具调用和审计归档等职责。

\subsection{D层核心动作}

D层围绕R层输出的三类结果执行动作。对于pass结果，D层保持原转账流程，完成必要的交易记录和审计写入。对于interrogate结果，D层向前端返回二次质询状态，提示用户补充说明或完成核验，再把用户回复交由I层和R层继续处理。对于block结果，D层阻断交易继续执行，向用户展示风险说明，并保存触发拦截的主要依据。

除自动化处置外，D层还提供人工复核入口。用户在风险报告页发起复核申请后，系统创建复核记录并进入后台处理队列。审核人员查看交易上下文、风险报告和审计日志后给出处置结论。复核结果作为治理记录保存，不直接覆盖原有自动化裁决，从而保证自动化链路和人工链路都可以被独立追踪。

\subsection{MCP协议与安全数据交互}

D层涉及交易执行、风险记录和审计落盘等关键动作，不能通过零散接口随意调用。本文使用 MCP 协议对这些能力进行封装，使上层只能通过明确的工具入口访问底层能力。MCP 在本系统中主要承担三项职责：一是\textbf{原子工具封装}，将风险计算、转账执行、日志写入等能力拆成边界清晰的工具；二是\textbf{权限管控}，每次调用都需要经过身份和范围校验，避免 Agent 接触未授权数据；三是\textbf{幂等校验}，以全局唯一交易流水号作为幂等键，避免重复请求导致重复执行。协议交互的完整时序如图所示。

% ===== 手动插图占位（已禁用，避免编译报错）=====
% 原图路径：figures/ui/candidate_p31_1_Im4.png
% \begin{figure}
%     \centering
%     \includegraphics[width=0.92\linewidth]{figures/ui/candidate_p31_1_Im4.png}
%     \caption{基于MCP协议的数据交互时序图}
%     \label{fig:mcp-sequence}
% \end{figure}

\subsection{全链路审计设计}
全链路审计用于回答三个问题：系统看到了什么、为什么做出当前判断、最终执行了什么动作。每笔交易从H层感知采集到D层指令执行，各处理节点都会生成对应审计事件。整个链路以全局唯一 Transfer\_ID 贯穿，审计事件链结构如图所示。

% ===== 手动插图占位（已禁用，避免编译报错）=====
% 原图路径：figures/ui/candidate_p31_2_Im5.png
% \begin{figure}
%     \centering
%     \includegraphics[width=0.82\linewidth]{figures/ui/candidate_p31_2_Im5.png}
%     \caption{全链路追踪流程图}
%     \label{fig:audit-flow}
% \end{figure}


审计日志记录感知快照、Agent 分析结果、R层评分与阈值路由、D层执行动作和人工复核状态等信息。后续审核人员可以通过交易流水号回溯各节点输入和输出，判断系统是否存在误判、漏判或状态流转异常。对于本文原型系统而言，审计链路不仅是合规要求，也是验证 HIRD 分层设计是否真正解耦的重要依据。

\section{本章小结}

本章围绕手机银行转账场景，对系统需求和总体架构进行了设计。3.1节从功能性和非功能性两个角度梳理了系统需要解决的问题，包括多模态感知、语义识别、权限约束、分级治理和审计追踪。3.2节提出 HIRD 四层架构，并说明感知、认知、治理和执行之间的数据流转关系。3.3节设计了H层的五项微行为指标和感知快照机制，3.4节说明I层中四类 Agent 的协作方式，3.5节给出R层的动态权重和分级路由方法，3.6节进一步说明D层如何完成业务执行、MCP 安全交互和全链路审计。

通过上述设计，本文把 AI Agent 的语义分析能力限制在认知层，把最终交易处置交给规则治理层和业务执行层完成。这样的分层方式既保留了大模型处理复杂文本的优势，也使交易执行、人工复核和审计记录保持确定性。第四章将在本章设计基础上说明系统的具体工程实现。

\chapter{系统实现与功能验证}

本章在第三章设计的基础上，说明原型系统的具体实现过程。实现工作主要围绕四个问题展开：移动端如何采集微行为数据，后端如何组织预检和二次核验流程，MCP 工具如何约束数据调用和交易执行，以及系统如何通过自动化测试验证主要功能是否闭环。与前文的架构设计相比，本章更关注代码层面的模块划分、接口流转和测试结果。

\section{系统开发环境与技术栈选型}

本文系统采用前后端分离的实现方式。移动端主要负责交易界面、行为采集和风险结果展示；后端负责风险预检、Agent 调度、规则治理、MCP 工具调用和审计记录。考虑到本文是本科毕业设计中的原型系统，开发环境以本地可复现实验为主，重点保证功能链路完整、接口可测试和数据可追踪。

\subsection{系统开发环境配置}

系统开发与测试环境如表\ref{tab:dev-env}所示。该环境能够支撑移动端界面调试、后端接口运行和自动化测试执行。

\begin{table}[H]
    \centering
    \caption{系统开发与运行环境配置表}
    \label{tab:dev-env}
    \begin{tabular}{|c|c|p{6cm}|}
        \hline
        环境类型 & 配置项 & 版本/规格说明 \\
        \hline
        \multirow{2}{*}{硬件环境} & CPU与内存 & 8核及以上处理器，16GB及以上内存 \\
        \cline{2-3}
        & 显卡 (GPU) & NVIDIA RTX 系列 \\
        \hline
        系统平台 & 操作系统 & win11 \\
        \hline
        \multirow{2}{*}{开发工具} & 后端 IDE & Visual Studio Code\\
        \cline{2-3}
        & 移动端 IDE & Flutter \\
        \hline
        \multirow{2}{*}{编程语言} & 后端语言 & Python 3.11\\
        \cline{2-3}
        & 移动端语言 & Dart 3.x\\
        \hline
    \end{tabular}
\end{table}

\subsection{核心技术栈选型}

围绕 HIRD 架构的实现需求，本文选择的核心技术栈如表\ref{tab:tech-stack}所示。各组件并不是简单堆叠，而是分别对应系统中的具体职责：Flutter 负责移动端交互和行为采集，FastAPI 负责后端异步接口，MCP 负责工具封装和权限边界，大语言模型与 FAISS 负责语义分析和案例检索，脱敏模块负责处理日志和模型输入中的敏感字段。

\begin{table}[H]
    \centering
    \caption{系统核心技术栈选型表}
    \label{tab:tech-stack}
    \begin{tabular}{|c|c|p{7cm}|}
        \hline
        架构层级 / 模块 & 技术组件 & 核心作用与选型优势 \\
        \hline
        感知集成层  & Flutter SDK & 跨平台高性能渲染，保证无感采集时 UI 流畅 \\
        \hline
        业务网关  & FastAPI & 异步非阻塞 IO，高并发下稳定可靠 \\
        \hline
        安全交互协议 & MCP Protocol & 标准化工具封装，实现安全解耦与权限隔离 \\
        \hline
        认知决策引擎 & 大语言模型 + FAISS & 检索增强生成，结合历史案例进行高维语义检索 \\
        \hline
        合规与审计模块 & PII 脱敏正则化 & 自动清洗敏感信息，满足金融隐私监管要求 \\
        \hline
    \end{tabular}
\end{table}

\section{基于 MCP 协议的安全数据中枢实现}

MCP 数据中枢的实现目标，是让 Agent 在分析交易风险时只能通过受控工具访问必要信息，而不是直接接触底层业务接口。本文将风险计算、转账执行、审计记录等能力封装为独立工具，并在工具内部加入权限校验、幂等检查和日志记录。这样做可以降低模型输出直接影响业务状态的风险，也便于后续排查每一次工具调用的来源和结果。

在后台 \texttt{mcp\_servers/bank\_server.py} 中，系统将核心业务能力封装为符合 MCP 契约的原子工具，例如风险等级判定工具和转账执行工具 \texttt{execute\_transfer}。这些工具以全局唯一交易流水号作为幂等键。当同一笔交易因为网络重试或并发请求触发重复调用时，工具会先检查流水号状态，再决定是否继续执行，从而避免重复扣款、重复审计或状态二次漂移。通过这种方式，Agent 只能够发起受控工具调用，不能绕过工具直接修改交易状态。

同时，系统在工具响应前调用 \texttt{server/app/services/pii\_masker.py} 中的脱敏处理逻辑，对银行卡号、手机号等敏感字段进行掩码转换。Agent 获取的是完成语义分析所需的脱敏上下文，而不是完整的个人身份信息。该设计与第三章提出的“分析权与执行权分离”保持一致：模型负责分析风险线索，交易执行和敏感数据处理仍由确定性代码控制。

\section{移动端无感微行为感知模块实现}

移动端感知模块对应第三章中的 H 层，主要负责在转账流程中采集用户操作过程数据。本文关注的是输入时长、非预期停顿、修改次数、页面停留和操作节奏方差等指标，而不是保存完整的用户输入过程。这样可以在获取风险判断线索的同时，减少对用户隐私和界面体验的影响\cite{li2024cognitive}。

实现上，系统利用 Flutter 的事件回调机制记录转账界面中的关键交互事件，包括字段输入、修改、页面停留和提交动作。为避免采集逻辑影响界面响应，相关统计尽量放在非阻塞流程中完成。用户提交交易时，移动端将行为统计结果与交易金额、收款信息、备注文本等字段一起封装为风险感知快照。

快照生成前，系统会对明显由误触、短时间重复点击或页面加载造成的异常事件进行过滤，再对连续型指标做归一化处理。处理后的快照通过后端接口进入预检流程。本文实现的重点是让 H 层输出格式稳定，便于后续 I 层、R 层和审计模块共同使用，而不是在前端直接给出风险结论。

\section{后端认知决策与异步编排实现}

后端实现围绕预检热路径和语义冷路径两条流程展开。预检热路径主要处理本地规则、结构化字段和状态记录，要求尽快返回；语义冷路径涉及 Agent、RAG 检索和大模型接口，耗时更长，需要异步挂起和状态机配合。FastAPI 在本文中承担接口编排和异步调度职责。

\textbf{1. 微服务协同与异步编排}

感知快照到达后端后，由 \texttt{server/app/api/routes/agent.py} 等路由接口接收并分发。系统通过 \texttt{llm\_gateway.py} 封装底层大模型调用，使上层逻辑不直接依赖某个具体模型接口；语义解析、案例检索和报告生成由 \texttt{agent\_service.py} 组织；最终评分和状态流转则由 \texttt{risk\_engine.py} 等模块处理。各模块之间通过交易流水号和审计事件保持关联。

\textbf{2. 双路径决策与冷热分流时延治理}

系统优先执行基于规则的热路径预检。第五章测试结果显示，该路径平均响应时间为 \textbf{43.60ms}，能够满足 100ms 内完成预检的设计目标。若交易进入二次核验或需要 Agent 语义分析，后端会将交易状态挂起，并通过异步任务等待模型返回。二次核验路径平均响应时间为 \textbf{2360.31ms}，大模型调用路径平均响应时间为 \textbf{5158.52ms}。由于冷路径耗时明显长于预检路径，前端需要通过等待提示和状态刷新承接这一延迟。

在结果收敛阶段，\texttt{risk\_engine.py} 会把 Agent 输出的语义线索、规则评分和动态权重结合起来，生成明确的 pass、interrogate 或 block 结果。这样做的目的，是让模型只参与分析，不直接输出业务动作。

\textbf{3. 全局时间戳与 Trace 审计}

系统在每一笔交易的生命周期内记录关键审计事件，包括预检开始、感知快照生成、Agent 分析、R层评分、D层执行和人工复核状态变化。每个事件绑定交易流水号和时间戳，后续可以按照时间线复盘处理过程。该机制主要用于定位问题和支持复核，而不是只在最终结果处记录一个简单状态。

\section{系统功能测试与验收验证}

为验证系统主要功能是否能够闭环，本文使用 PyTest 自动化测试框架对后端接口、风险公式、二次核验、人工复核和 MCP 审计链路进行验证。测试数据以场景化构造样本为主，包括常规转账、虚假理财、冒充公检法、冒充客服和异常停顿等类型。第四章只对测试结果做工程验收说明，完整实验统计放在第五章展开。

\textbf{1. 核心算法与自动化接口验收}

通过执行 \texttt{server/tests/test\_api.py} 与 \texttt{test\_formula.py}，系统验证了接口返回、特征归一化和动态权重计算等基础逻辑。第五章统计显示，全部 \textbf{66} 个后端测试用例均通过 PyTest 自动化验证；在 \textbf{54} 个核心场景样本中，系统精确匹配率为 \textbf{92.6\%}，欺诈召回率为 \textbf{100\%}，正常交易误报率为 \textbf{0\%}。这些结果说明，当前原型系统在构造样本范围内能够完成预期判断，但仍需要在更大规模真实数据上继续验证泛化能力。

\textbf{2. 柔性管控与性能时延验收}

针对双路径时延和状态机流转，本文执行 \texttt{server/tests/test\_manual\_review.py} 等测试脚本。测试结果显示，预检热路径平均响应时间为 \textbf{43.60ms}，P95 为 \textbf{78.39ms}；二次核验路径平均响应时间为 \textbf{2360.31ms}，P95 为 \textbf{2767.42ms}；大模型调用路径平均响应时间为 \textbf{5158.52ms}。这些数据与系统设计预期一致：常规交易走快速路径，复杂交易通过异步状态机承接较长的模型调用。

\textbf{3. MCP 审计合规性溯源验收}

通过 \texttt{test\_trace\_mcp.py}，系统对 Trace ID 传递、MCP 工具幂等性和日志脱敏效果进行了验证。测试日志显示，并发重放请求会被工具状态校验逻辑拦截；落盘日志经过 \texttt{pii\_masker.py} 处理后，银行卡号、手机号等字段能够被掩码替换。该结果说明，MCP 工具链和审计模块在当前测试范围内能够满足原型系统的可追踪要求。

\section{移动端界面展示与核心交互流程}

本节展示移动端主要界面和风控交互流程。前端界面除了承载常规转账操作，还需要展示二次质询、风险报告、AI助手说明和人工复核入口。与第三章中的D层设计对应，前端不重新判断风险，只负责把后端返回的状态和说明呈现给用户，并把用户补充信息提交回后端。

\subsection{首页与常规转账交互}

系统首页采用接近手机银行的常规布局，顶部展示账户资产信息，下方提供账单查询和转账入口。为了方便测试不同风险流程，页面中额外设置了风险演示入口，用于触发预设场景。

在低风险交易场景下，用户填写转账信息并提交后，后端走预检热路径完成判定，前端根据 pass 结果继续交易流程。由于此类交易不需要二次质询，用户侧不会看到额外风控弹窗。

% ===== 手动插图占位（已禁用，避免编译报错）=====
% 原图路径：figures/ui/candidate_p37_1_Im6.png
% \begin{figure}
%     \centering
%     \includegraphics[width=0.92\linewidth]{figures/ui/candidate_p37_1_Im6.png}
%     \caption{首页}
%     \label{fig:ui-home}
% \end{figure}
该界面承载 H 层输入上下文（账户、入口、城市与基础交易环境），用于在不打断主流程的前提下形成首轮风险感知快照。

% ===== 手动插图占位（已禁用，避免编译报错）=====
% 原图路径：figures/ui/candidate_p38_1_Im7.png
% \begin{figure}
%     \centering
%     \includegraphics[width=0.92\linewidth]{figures/ui/candidate_p38_1_Im7.png}
%     \caption{转账窗口}
%     \label{fig:ui-transfer}
% \end{figure}
转账窗口对应主流程“提交前采集”阶段，金额、收款方与设备上下文在此汇总后进入预检与治理链路。


\subsection{二次质询交互流程}

当后端预检结果为 interrogate 时，前端弹出二次质询窗口，要求用户补充说明转账背景。该窗口并不直接判断用户回答是否安全，而是把用户输入作为新的语义分析材料，提交给后端继续处理。

弹窗问题主要围绕三类信息展开：用户与收款方关系、转账用途、操作环境是否涉及验证码索要、安全账户或屏幕共享等风险点。用户提交说明后，前端调用 \texttt{/api/v1/transfers/secondary-check} 接口将文本发送至后端。后端先执行本地红旗规则，再在必要时调用 Agent 语义分析，最后由治理层收敛为新的处置结果。

若二次质询通过，交易继续执行；若回复命中高风险规则或语义分析结果指向欺诈，前端展示拦截提示，并提供风险报告入口。该流程用于处理中风险交易，避免系统在证据不足时直接拦截。

% ===== 手动插图占位（已禁用，避免编译报错）=====
% 原图路径：figures/ui/candidate_p38_2_Im8.png
% \begin{figure}
%     \centering
%     \includegraphics[width=0.92\linewidth]{figures/ui/candidate_p38_2_Im8.png}
%     \caption{二次质询窗口}
%     \label{fig:ui-interrogate}
% \end{figure}
二次质询页承接 R 层的 interrogate 指令，用户补充说明作为 I 层语义复判输入，并决定后续放行或拦截。


\subsection{风险报告查看}

系统在“我的”页面中设置风险报告入口，用于展示被拦截或进入风险处置流程的交易记录。用户点击记录后，可以查看该笔交易的结构化风险报告。

报告面板主要展示风险标题（headline）、综合风险等级（overall\_risk\_level）、风险摘要（risk\_summary）和风险因子列表（risk\_factors）。同时，报告还保留报告版本、策略版本、生成阶段和关联审计事件等治理元数据。这样用户和审核人员可以看到系统给出当前处置的主要依据，而不是只看到一个最终状态。

% ===== 手动插图占位（已禁用，避免编译报错）=====
% 原图路径：figures/ui/candidate_p39_1_Im10.png
% \begin{figure}
%     \centering
%     \includegraphics[width=0.92\linewidth]{figures/ui/candidate_p39_1_Im10.png}
%     \caption{风险报告}
%     \label{fig:ui-risk-report}
% \end{figure}
风险报告是“处置与审计”界面，集中展示风险证据、治理上下文与关联事件，支撑复核申请与审计留痕。


\subsection{AI助手与人工复核}

系统提供 AI 助手入口，用于解释风险报告中的主要内容。用户可以询问交易被拦截的原因，助手会基于该笔交易的报告摘要和审计记录生成解释。该功能只用于说明，不改变交易状态，也不替代人工复核。

对于被拦截的交易，用户可以在报告详情页发起人工复核。复核申请提交后，交易进入后台处理队列。审核人员可在复核处理台查看待处理记录，并按 submitted、in\_review、closed 等状态筛选。接单后，审核人员可以查看交易上下文、风险报告和审计链路，再给出处置结论。复核详情页按时间线展示 submitted\_at、in\_review\_at 和 closed\_at 三个节点，复核结果作为治理记录留存，不直接回写自动化裁决。

% ===== 手动插图占位（已禁用，避免编译报错）=====
% 原图路径：figures/ui/candidate_p39_2_Im9.png
% \begin{figure}
%     \centering
%     \includegraphics[width=0.92\linewidth]{figures/ui/candidate_p39_2_Im9.png}
%     \caption{AI助手}
%     \label{fig:ui-ai-assistant}
% \end{figure}
AI助手用于解释报告结论和核验原因，定位为解释与沟通组件，不直接修改 D 层最终执行状态。


上述界面共同构成系统执行层的前端流程：低风险交易直接进入常规转账，中风险交易触发二次质询，高风险交易展示拦截和报告入口，存在争议的交易可以发起人工复核。每一步交互都对应后端接口和审计事件，便于后续复盘。

\section{本章小结}

本章围绕原型系统的工程实现进行了说明。首先介绍了系统开发环境和技术栈，并说明 Flutter、FastAPI、MCP、大语言模型、FAISS 和脱敏模块分别承担的职责。随后，说明了 MCP 数据中枢、移动端微行为采集、后端双路径编排和审计事件记录的实现方式。最后，结合自动化测试结果，对接口、公式、状态机、MCP 工具和日志脱敏等功能进行了验收说明。

从当前测试结果看，系统全部 66 个后端测试用例能够通过自动化验证；在 54 个核心场景样本中，精确匹配率为 92.6\%，欺诈召回率为 100\%，正常交易误报率为 0\%。性能方面，预检热路径平均响应时间为 43.60ms，二次核验路径平均响应时间为 2360.31ms，大模型调用路径平均响应时间为 5158.52ms。上述结果说明，本文原型系统在构造测试集范围内完成了主要业务闭环。第五章将进一步对这些测试数据进行集中分析。

\chapter{实验测试与分析}

本章对 HIRD 智能风控原型系统进行测试与分析。测试重点包括四部分：一是后端各模块的功能测试，确认接口、算法、MCP 工具和审计模块能够按预期运行；二是三条主要路径的响应时延测试，观察预检、二次核验和大模型调用在本地环境下的耗时差异；三是基于 54 个场景样本的识别结果对比；四是对 H 层微行为特征和 R 层动态权重机制进行消融观察。需要说明的是，本章测试数据主要来自构造样本和自动化脚本，实验结果用于验证原型系统在当前测试条件下的表现，不等同于真实银行生产环境中的最终效果。

\section{实验环境与数据集配置}

实验在本地开发环境中执行，硬件配置为 8 核处理器、16GB 内存，操作系统为 Windows 11，Python 版本为 3.11，后端框架为 FastAPI，测试框架为 PyTest。系统后端服务以单进程方式运行，未引入额外负载均衡、容器编排或分布式压测环境。因此，本章性能结果主要反映轻载条件下的接口响应情况。

实验数据分为两类。第一类是功能测试用例集，覆盖感知快照生成、MCP 工具调用、认知层 Agent 推理、治理层公式计算、执行层指令下发和审计日志落盘等模块，共计 66 个 PyTest 自动化测试用例。第二类是场景化精度测试集，共包含 54 个经过标注的样本，覆盖正常转账、语义欺诈变种和边界模糊案例，用于观察系统在不同风险等级下的判定结果。

\section{系统功能测试}

功能测试用于检查系统主要模块是否能够按照设计完成输入处理、状态流转和结果输出。本文基于 PyTest 对后端模块进行回归测试，测试范围包括核心算法模块、MCP 协议模块、认知决策模块、感知集成模块，以及执行与审计模块。测试结果如表\ref{tab:func-test-results}所示。

\begin{table}[H]
    \centering
    \caption{系统功能测试结果汇总}
    \label{tab:func-test-results}
    \begin{tabular}{|c|c|c|c|c|}
        \hline
        测试模块 & 用例数量 & 通过数量 & 通过率 & 覆盖内容 \\
        \hline
        核心算法模块（\texttt{test\_formula.py}） & 15 & 15 & 100\% & 风险评分计算、动态权重调节、三级阈值路由 \\
        \hline
        MCP协议模块（\texttt{test\_trace\_mcp.py}） & 12 & 12 & 100\% & 原子工具封装、幂等校验、PII脱敏 \\
        \hline
        认知决策模块（\texttt{test\_agent.py}） & 14 & 14 & 100\% & Agent协作流程、语义分析、RAG检索 \\
        \hline
        感知集成模块（\texttt{test\_harvest.py}） & 11 & 11 & 100\% & 快照生成、特征归一化、噪声过滤 \\
        \hline
        执行与审计模块（\texttt{test\_delivery.py}） & 14 & 14 & 100\% & 指令执行、审计日志落盘、TraceID穿透 \\
        \hline
        合计 & 66 & 66 & 100\% & — \\
        \hline
    \end{tabular}
\end{table}

从测试结果看，66 个用例均通过断言验证，说明当前代码版本中的主要后端功能链路能够正常运行。其中，核心算法模块主要验证风险评分和阈值路由是否按配置输出；MCP 协议模块重点验证工具封装、幂等校验和脱敏逻辑；执行与审计模块则用于确认交易状态变化后能否正确写入日志。该结果只能说明系统在既定测试用例覆盖范围内通过验证，尚不能代表所有真实业务场景。

\section{系统性能指标评估}

性能测试关注三条主要路径的端到端响应时间：预检热路径、二次核验路径和大模型调用路径。每条路径独立执行 100 次基准测试，记录平均响应时间、P50、P95、最大值和最小值。三条路径的处理复杂度不同，因此本节分别讨论。

\subsection{预检热路径时延}

预检路径是每笔交易都会经过的首轮风控判断。该路径只涉及本地规则匹配、结构化字段计算和必要的状态记录，不主动调用大模型。测试结果如表\ref{tab:precheck-latency}所示。

\begin{table}[H]
    \centering
    \caption{预检热路径时延统计}
    \label{tab:precheck-latency}
    \begin{tabular}{|c|c|}
        \hline
        指标 & 数值 \\
        \hline
        平均响应时间 & 43.60 ms \\
        \hline
        P50响应时间 & 41.20 ms \\
        \hline
        P95响应时间 & 78.39 ms \\
        \hline
        最大响应时间 & 89.50 ms \\
        \hline
        最小响应时间 & 38.10 ms \\
        \hline
    \end{tabular}
\end{table}

预检热路径平均响应时间为 43.60ms，P95 为 78.39ms，最大值为 89.50ms。在本地轻载测试条件下，该路径能够满足 NF-01 中“预检≤100ms”的约束。由于该路径不涉及外部模型请求，时延主要来自接口处理、本地规则计算和日志记录。

\subsection{二次核验路径时延}

当预检结果为 interrogate 时，系统会触发二次核验流程。该路径需要接收用户补充说明，并进行语义分析和状态更新。测试结果如表\ref{tab:secondary-check-latency}所示。

\begin{table}[H]
    \centering
    \caption{二次核验路径时延统计}
    \label{tab:secondary-check-latency}
    \begin{tabular}{|c|c|}
        \hline
        指标 & 数值 \\
        \hline
        平均响应时间 & 2360.31 ms \\
        \hline
        P50响应时间 & 2310.15 ms \\
        \hline
        P95响应时间 & 2767.42 ms \\
        \hline
        最大响应时间 & 2980.20 ms \\
        \hline
    \end{tabular}
\end{table}

二次核验路径平均时延为 2360.31ms，P95 为 2767.42ms，最大值为 2980.20ms。该结果处于 3s 以内，符合本文对二次交互路径的设计约束。与预检路径相比，该路径耗时明显增加，主要原因是它包含用户补充说明后的语义分析和状态机流转。前端通过异步等待承接该耗时，因此不会阻塞界面主流程。

\subsection{大模型调用路径时延}

当本地规则无法覆盖二次核验内容时，系统需要调用公网大模型 API 进行进一步语义分析。该路径受网络往返、模型推理和返回内容长度影响较大，测试结果如表\ref{tab:llm-path-latency}所示。

\begin{table}[H]
    \centering
    \caption{大模型调用路径时延统计}
    \label{tab:llm-path-latency}
    \begin{tabular}{|c|c|}
        \hline
        指标 & 数值 \\
        \hline
        平均响应时间 & 5158.52 ms \\
        \hline
        P50响应时间 & 5020.40 ms \\
        \hline
        P95响应时间 & 5635.89 ms \\
        \hline
        最大响应时间 & 6120.70 ms \\
        \hline
    \end{tabular}
\end{table}

大模型调用路径平均时延为 5158.52ms，P95 为 5635.89ms。该路径耗时高于预检和二次核验路径，说明公网模型调用不适合放入所有交易的实时预检链路中。本文系统因此采用冷热分流策略：常规交易优先走本地预检路径，只有在需要深层语义分析时才进入大模型路径，并通过协程调度和前端异步提示维持流程连续性。

\section{对比实验}

本节通过场景化样本测试观察 HIRD 架构在当前测试集上的识别结果，并通过影子模式检查预检路径输出是否稳定。由于样本规模有限，本节结果主要用于说明原型系统在构造样本中的表现，不能直接推导到所有真实交易场景。

\subsection{识别精度对比}

本文使用 54 个场景化测试样本进行端到端精度评估，样本包含欺诈变种、正常交易和边界模糊案例。测试结果如表\ref{tab:scenario-accuracy}所示。

\begin{table}[H]
    \centering
    \caption{场景化精度测试结果}
    \label{tab:scenario-accuracy}
    \begin{tabular}{|c|c|c|}
        \hline
        指标 & 数值 & 说明 \\
        \hline
        样本总数 & 54 & 含欺诈变种与正常交易 \\
        \hline
        精确匹配数 & 50 & 判定结果与标注完全一致 \\
        \hline
        精确匹配率 & 92.6\% & 50/54 \\
        \hline
        欺诈召回率 & 100\% & 所有欺诈样本均被正确识别 \\
        \hline
        正常交易误报率（block-only） & 0\% & 所有正常样本均未被误拦截 \\
        \hline
    \end{tabular}
\end{table}

在 54 个样本中，系统有 50 个样本的判定结果与标注完全一致，精确匹配率为 92.6\%；欺诈召回率为 100\%，说明测试集中的欺诈样本均被识别；正常交易误报率（block-only）为 0\%，说明正常样本没有被直接误拦截。需要注意的是，当前样本属于结构化可控测试集，欺诈文本特征相对明确，尚未覆盖真实生产环境中的长尾话术、复杂用户习惯和高并发干扰。因此，该结果更适合作为原型验证依据，而不是生产级泛化结论。

\subsection{预检路径数据一致性验证}

为检查预检路径在相同输入下是否稳定，本实验引入影子模式测试。主路径和影子路径执行相同逻辑，仅在调用链路上进行镜像观测，然后比较两者输出是否一致。测试结果如表\ref{tab:shadow-mode-test}所示。

\begin{table}[H]
    \centering
    \caption{预检路径影子模式测试结果}
    \label{tab:shadow-mode-test}
    \begin{tabular}{|c|c|}
        \hline
        指标 & 数值 \\
        \hline
        总比对次数 & 100 \\
        \hline
        偏离次数 & 0 \\
        \hline
        偏离率（drift rate） & 0.0\% \\
        \hline
    \end{tabular}
\end{table}

影子模式测试中，总比对次数为 100 次，偏离次数为 0，偏离率为 0.0\%。这说明在本次测试输入和代码版本下，预检主路径与影子路径能够输出一致结果。该测试主要验证规则路径的确定性，不涉及大模型输出波动。

\section{消融实验}

为观察 HIRD 架构中关键模块对当前测试集的影响，本节进行了两组消融实验：一组移除 H 层微行为特征输入，另一组将 R 层动态权重调节因子 $\omega_{adj}$ 设置为 0。实验仍使用 54 个场景样本，观察精度指标变化。

\textbf{H层感知特征消融}：移除H层采集的微行为特征（输入时长、停顿频率、操作节奏方差等五项指标），仅依赖交易金额、设备IP等结构化数据与交易附言文本进行风险判定。

\textbf{R层动态权重消融}：将公式（3-1）中的动态权重调节因子 $\omega_{adj}$ 强制设置为0，模拟移除R层的决策收敛机制，允许规则评分直接决定最终判定。

测试结果如表\ref{tab:ablation-study}所示。

\begin{table}[H]
    \centering
    \caption{消融实验结果对比}
    \label{tab:ablation-study}
    \begin{tabular}{|c|c|c|c|c|}
        \hline
        实验组 & 欺诈召回率 & 正常交易误报率（block-only） & 精确匹配率 & 说明 \\
        \hline
        完整HIRD系统 & 100\% & 0\% & 92.6\% & 全模块运行 \\
        \hline
        H层消融（无微行为特征） & 100\% & 0\% & 92.6\% & 结构性特征+文本仍有效 \\
        \hline
        R层消融（$\omega_{adj}=0$） & 100\% & 0\% & 92.6\% & $\omega_{adj}$在此测试集上影响有限 \\
        \hline
    \end{tabular}
\end{table}

从表中可以看到，在当前 54 个场景样本上，移除 H 层微行为特征或关闭动态权重调节后，三个主要指标没有发生变化。该结果说明，本测试集中不少欺诈样本已经包含较明显的文本或结构化风险特征，系统即使不依赖微行为特征，也能够完成识别。换言之，本组消融实验没有充分体现 H 层和动态权重机制的独立贡献。

这一结果并不意味着 H 层和 R 层设计没有价值，而是提示当前测试集还不够“难”。如果后续要进一步验证微行为特征的作用，需要增加弱文本特征样本，例如交易备注较正常、但用户操作节奏异常的场景；如果要验证动态权重机制，需要增加规则分数与语义分数冲突的边界样本。该局限将在第六章中继续讨论。

\section{本章小结}

本章基于本地测试环境，对 HIRD 智能风控原型系统进行了功能、性能、场景精度、一致性和消融测试。功能测试中，66 个后端测试用例均通过 PyTest 自动化验证。性能测试中，预检热路径平均响应时间为 43.60ms，P95 为 78.39ms；二次核验路径平均响应时间为 2360.31ms，P95 为 2767.42ms；大模型调用路径平均响应时间为 5158.52ms。上述结果说明，不同路径的耗时差异明显，也印证了冷热分流设计的必要性。

场景精度测试中，54 个样本的精确匹配率为 92.6\%，欺诈召回率为 100\%，正常交易误报率为 0\%。影子模式测试中，预检路径偏离率为 0.0\%，说明本地规则路径在相同输入下具有稳定输出。消融实验中，H 层与 R 层消融后指标未发生变化，说明当前样本仍偏向显性文本和结构化风险特征，对微行为特征和动态权重机制的独立贡献验证不足。综合来看，本章测试能够证明原型系统在当前构造样本和本地环境下具备可运行性和基本闭环能力，但仍需要在更复杂样本和更接近真实业务的环境中继续验证。

\chapter{结论与展望}

本文围绕手机银行转账场景中的语义欺诈防控问题，完成了从需求分析、架构设计、核心模块实现到自动化测试验证的完整研究过程。与传统依赖金额、地域、设备状态等结构化字段的风控方式相比，本文更关注用户在被诱导或胁迫状态下产生的操作节奏变化，并将微行为特征、交易文本语义和确定性规则治理结合起来，形成一套基于 AI Agent 与 MCP 协议的 HIRD 分层风控原型系统。需要说明的是，本文实验主要依托模拟场景数据和项目自动化测试展开，相关结论主要成立于本文构造的测试条件之内。

\section{主要研究结论}

结合系统实现和第五章测试结果，本文得到以下结论：

1、\textbf{HIRD 分层架构能够支撑手机银行转账场景下的风险闭环。} 本文将系统划分为感知集成层、认知决策层、规则治理层和业务执行层，分别承担数据采集、语义分析、决策收敛和动作执行职责。在工程实现中，移动端负责生成微行为快照，后端负责预检、二次核验、风险报告和人工复核等流程编排。第五章测试结果显示，系统全部 66 个后端测试用例均通过 PyTest 自动化验证，通过率为 100\%，说明各模块在当前实现范围内能够完成预期业务闭环。

2、\textbf{双路径决策机制在实时性和深度分析之间取得了可验证的平衡。} 对于常规交易，系统优先执行预检热路径，由本地规则和结构化字段完成快速判定；对于需要进一步确认的中高风险交易，系统再进入二次核验或 Agent 语义分析路径。性能测试表明，预检热路径平均响应时间为 43.60ms，P95 为 78.39ms，满足 100ms 内完成预检的设计目标；二次核验路径平均响应时间为 2360.31ms，P95 为 2767.42ms，满足异步交互场景下 3s 内响应的约束。大模型调用路径平均耗时为 5158.52ms，虽然明显高于本地规则路径，但通过异步挂起和状态机流转，系统仍能保持业务流程的连续性。

3、\textbf{场景化测试验证了系统对典型欺诈样本的识别能力。} 本文基于 54 个核心场景样本，对常规转账、虚假理财、冒充公检法、冒充客服和异常停顿等场景进行测试。实验结果显示，系统精确匹配率为 92.6\%，欺诈召回率为 100\%，正常交易误报率为 0\%。这些结果说明，在当前构造的测试集中，系统能够较稳定地区分正常交易与典型语义欺诈交易。同时，影子模式测试中预检路径判定偏离率为 0.0\%，表明在相同输入条件下，系统主路径与镜像路径能够保持一致输出。

4、\textbf{MCP 协议为工具调用和审计记录提供了统一边界。} 系统将风险计算、转账执行、审计记录等能力封装为 MCP 原子工具，并通过 Trace ID 串联预检、Agent 分析、规则治理和执行结果。该设计使 Agent 主要承担语义线索分析，不直接操作资金执行动作；最终业务动作仍由确定性代码逻辑接管。通过这种方式，系统在利用大模型语义理解能力的同时，为后续审计、复核和问题定位保留了较清晰的调用链路。

\section{本文创新点总结}

本文工作主要体现在以下三个方面：

1、\textbf{将微行为快照引入手机银行转账风控链路。} 本文没有只依赖金额、设备、IP 等静态字段，而是从用户操作过程入手，将输入时长、停顿频率、修改次数、页面停留和操作节奏方差等特征封装为风险快照。该设计为识别“交易参数正常但用户可能被诱导”的场景提供了额外观察维度。

2、\textbf{设计了分析权与执行权分离的 HIRD 架构。} 本文将 AI Agent 放置在认知决策层，主要负责文本语义解析、风险线索提取和报告生成；规则治理层负责将语义结果与硬规则评分进行整合；业务执行层只处理放行、二次核验、拦截和人工复核等确定性动作。这种结构降低了 Agent 直接影响资金流转的风险，也使系统边界更容易被测试和审计。

3、\textbf{在原型系统中落地了 MCP 风控工具链。} 系统围绕 MCP 协议封装风险计算、转账执行和审计记录等工具，并结合幂等校验、状态机流转和 Trace ID 追踪机制，形成了一条可复查的调用链路。相比普通接口调用，该方式更适合需要权限边界、日志追踪和复核依据的智能风控场景。

\section{研究不足与未来展望}

本文虽然完成了原型系统开发和自动化测试验证，但仍存在以下不足。

首先，当前实验主要依托 54 个核心场景样本和 2000 条场景化模拟数据展开，尚未接入真实银行脱敏流水。模拟数据能够覆盖典型诈骗话术和异常操作节奏，但难以完整反映真实业务中的长尾用户习惯、复杂交易关系和突发并发压力。因此，本文得到的精确匹配率、召回率、误报率和时延数据，仍需要在更大规模、更贴近真实业务的数据集上继续验证。

其次，第五章消融实验显示，在当前 54 个样本中，移除 H 层微行为特征或将 $\omega_{adj}$ 设置为 0 后，系统精度指标未出现明显下降。该现象说明当前样本中的欺诈文本特征较为显性，结构化规则和语义分析已经能够覆盖大部分风险判断。换言之，本文尚未充分构造出能够体现微行为特征和动态权重机制独立贡献的隐蔽样本，这也是后续实验设计需要重点补强的部分。

再次，系统目前主要围绕手机银行转账业务展开，尚未覆盖理财申购、贷款审批、账户变更等更多金融操作。不同业务的风险目标、用户行为路径和合规要求并不相同，现有 HIRD 架构虽然具备扩展空间，但 MCP 工具定义、阈值设置和前端交互方式仍需要结合具体业务重新设计。

后续研究可从三个方向继续推进：第一，扩展真实或半真实的风险案例库，增加弱文本特征、长尾诈骗话术和复杂微行为模式样本，使多模态感知与动态权重机制的贡献能够被更充分地验证；第二，结合更多用户历史行为与设备环境特征，优化 Agent 语义分析与规则治理之间的协同方式，降低对显性关键词的依赖；第三，在更接近生产环境的并发条件下继续测试 MCP 工具链的权限控制、幂等防护、异常恢复和审计追踪能力，使系统逐步从课程原型走向更接近真实金融业务的工程形态。

\end{document}
